%% sec_background.tex
%% ─────────────────────────────────────────────────────────────────────────────
%% Section 2 — Background and Related Work
%% ─────────────────────────────────────────────────────────────────────────────

\section{Background and Related Work}\label{sec:background}

\subsection{MC event generation and the matrix element bottleneck}

Physics event generators model proton collisions by integrating the transition
probability
\begin{equation}
  \mathcal{P}(i\to\mathcal{X})
  = \left|\mathcal{M}_{i\to\mathcal{X}}\right|^2,
  \label{eq:me}
\end{equation}
over phase space using MC sampling.  The calculation is \emph{embarrassingly
parallel}: the same \ME evaluation is performed independently for each sampled
PSP.  Modern generators such as \amcnlo provide automatic generation of the
amplitude code from Feynman rules via the HELAS
library~\cite{Murayama:1992gi}, where each diagram corresponds to a sequence
of wavefunction and vertex calls accumulating into colour-flow amplitude sums
(\jamp).

\subsection{\ggttg as validation benchmark}\label{sec:process}

We target the \ggttg process at leading order.  \mgfive generates 16 Feynman
diagrams ($\mathrm{D1}$--$\mathrm{D16}$, where $\mathrm{D16}$ has three
sub-diagrams via four-gluon contact vertices), 6 colour-flow amplitudes
($\jamp[0..5]$), and 32 helicity configurations.  The complete HELAS call
graph involves 10 distinct functions:
\texttt{ixxxxx}, \texttt{oxxxxx}, \texttt{vxxxxx} (external wavefunctions);
\texttt{FFV1\_0/1/2}, \texttt{FFV1P0\_3} (fermion--fermion--vector vertices);
\texttt{VVV1P0\_1}, \texttt{VVV1\_0} (triple-gluon vertices);
\texttt{VVVV1/3/4P0\_1} (four-gluon contact vertices).

Although the process is relatively compact (5 external particles), it provides
a demanding PM stress-test: the monolithic kernel implementing all 16 diagrams
requires ${\sim}38\KB$ of program memory---$2.4\times$ the 16\KB tile limit.
This makes \ggttg the minimal realistic benchmark for the partitioning
methodology developed here.

\subsection{GPU and vectorised-CPU acceleration}

The CUDACPP plugin for \mgfive~\cite{Valassi:2025xfn} achieves
$100\times$--$1000\times$ speedup over single-threaded Fortran on NVIDIA V100
GPUs for processes such as $\gluglu\to\ttbar\gluglu\gluglu$.  CPU vector
implementations (AVX2/AVX-512) deliver up to $16\times$ on a single core.
These results establish the performance target against which the AIE
architecture is compared in Section~\ref{sec:results}.

\subsection{Prior FPGA acceleration work}

Barbone et al.~\cite{Carrazza:2021gpx} demonstrated ${\sim}100\times$ CPU
speedup for $e^+e^-\to\mu^+\mu^-$ on a Xilinx Alveo using HLS-generated
circuits.  Barbone et al.~\cite{Barbone:2023mul} subsequently extended this to
the Versal platform, combining PL HLS kernels, VHDL datapaths, and \aie
kernels for the same two-body process.  Our work differs in two key respects:
(i)~we target a significantly more complex hadronic process (\ggttg with 16
diagrams vs.\ 1--2 for $e^+e^-\to\mu^+\mu^-$), and (ii)~we identify PM
constraints as the primary architectural driver, motivating a multi-tile
cascade pipeline rather than a single-core solution.
